\section{Characterizing Long-Context dLLM Inference}
\label{sec:char}

\subsection{A Per-Step Cost Model}
With a prefix KV cache, one denoising step processes a block of $B$ query tokens against a context of
$L$ tokens on $G$ GPUs. For MoE models we model the step time as
\begin{equation}
\begin{split}
T_{\text{step}} \approx\;&
\underbrace{\frac{S_{\text{KV}}(L)}{G\cdot BW_{\text{HBM}}}}_{\text{KV read}} +
\underbrace{T_{\text{exp}}(B)}_{\text{expert weights}} +
\underbrace{T_{\text{EP}}(B)}_{\text{dispatch/combine}} \\
&+ \underbrace{T_{\text{CP}}(B, G)}_{\text{attention merge}} + T_{\text{samp}},
\end{split}
\label{eq:step}
\end{equation}
where $S_{\text{KV}}(L) = 2\,n_{\text{layers}}\,n_{\text{kv}}\,d_{\text{head}}\,b\,L$ bytes. The
decode throughput of one request is $\TPF / T_{\text{step}}$. An AR model pays a comparable
$T_{\text{step}}$ per \emph{single} token, so when the KV-read term dominates, the dLLM advantage
approaches \TPF. Prefill costs $\Theta(L^2)$ attention FLOPs; fully bidirectional prompts cannot use
the causal $2\times$ saving that block-causal prompts retain.

\begin{table}[t]
\caption{KV-cache footprint from model configurations (BF16). Block-diffusion LLaDA2.x models and
their AR counterparts (Ling 2.0) share identical shapes.}
\label{tab:kv}
\centering\small
\resizebox{\columnwidth}{!}{%
\begin{tabular}{lrrr}
\toprule
Model & KV heads & KiB/token & GB @ 2M \\
\midrule
LLaDA2.2-mini / Ling-mini-2.0 (16B, 1.4B act.) & 4 & 40 & 82 \\
LLaDA2.2-flash / Ling-flash-2.0 (103B, 6.1B act.) & 4 & 64 & 131 \\
SDAR-30B-A3B & 4 & 96 & 197 \\
Qwen2.5-7B-Instruct-1M (AR) & 4 & 56 & 115 \\
LLaDA-MoE-7B (bidirectional) & 16 & 128 & 262 \\
LLaDA-8B (bidirectional) & 32 & 512 & 1{,}049 \\
\bottomrule
\end{tabular}}
\end{table}

\subsection{Where the Time Goes}
\placeholderfig{fig:breakdown}{Step-time breakdown versus context length (4K--2M) for LLaDA2.2-mini
and LLaDA2.2-flash on 8 MI355X GPUs.}{F1: stacked bars per context length: KV read, expert weights,
EP dispatch/combine, CP merge, sampling. Measured with the profiling pipeline of
\S\ref{sec:impl-prof}.}

\todo{Describe F1 once measured.} Our analytical estimate for LLaDA2.2-flash at 2M tokens on one node
(batch 1) is \est{$\approx$7.5\,ms} per step, split roughly evenly between KV reads, expert-weight
reads and EP communication, which motivates attacking all three.

\subsection{Diffusion versus Autoregressive at Equal Architecture}
\placeholderfig{fig:arvsdllm}{Decode throughput versus context length: LLaDA2.2 versus Ling 2.0
(identical architecture), and Qwen2.5-1M.}{F2: tokens/s per request (log scale) vs context 32K--2M.
Expect a crossover where the dLLM advantage grows toward \TPF.}

\hyp{The dLLM advantage grows with context and approaches the measured \TPF (2--5).}

\subsection{Gaps in Existing Serving Systems}
\label{sec:char-gaps}
SGLang serves LLaDA2.x, SDAR and DiffusionGemma, but its dLLM path currently disables
prefill/decode disaggregation, the hierarchical KV cache (and therefore storage backends such as
UMBP), pipeline parallelism, and, on AMD GPUs, graph capture~\cite{sglangdllm}. Context parallelism
is not disabled but is untested with dLLMs. vLLM serves DiffusionGemma natively and LLaDA2 only via
an experimental single-request plugin~\cite{diffusiongemma}. dInfer~\cite{ma2025dinfer} and
DiLaServe~\cite{chang2026dilaserve} target a single node or tensor parallelism only.
